\paragraph{Problem 3 answer} \GradescopeComment{This should be on page 10 (top) of your PDF.}

Consider any sequence $(\ell, \ell+1, \ldots, u)$ of more than two stops,
with middle stop $m=\lfloor (\ell+u)/2\rfloor$.

\begin{enumerate}[label=(\alph*)]

\item % (2a)
  Define $Q(\ell, u)$ as follows:
  \begin{blue}
    \[
      Q(\ell, u)
      =
      \REPLACEME
    \]

    % USEFUL LATEX MACROS (USE IN MATH MODE):
    % 
    % Use \segment{X}{Y} to represent a segment from X to Y.
    % 
    % Use A \cup B for set union.
    % 
    % Use \{ .. \} to make curly braces to define a set.  For example
    % 
    % \{ \segment i j : 1\le i < j \le N \} 
    % 
    % is the set of all $N \choose 2$ possible segments.
  \end{blue}

  % [review 2026-09-27] fix: number of segments is n-1 = u-l, not l-u
  Then $Q(\ell, u)$ has at most $n-1 = u-\ell$ segments,
  % [review 2026-09-27] fix: j is in the second half of the stops, so m < j (was m <= j), as in the problem statement
  and for any $(i, j)$ with $\ell \le i \le m < j \le u$
  the hop length from $i$ to $j$ in $Q(\ell, u)$ is at most two.

\item % (2b)
  Define $S(\ell, u)$ as follows:
  \BLUE{
    \[
      S(\ell, u) = \REPLACEME
    \]
  }
\end{enumerate}

\continued 
\setlength{\ITEMHT}{1in}

\paragraph{Problem 3(c) answer} \GradescopeComment{This should be on page 11 (top) of your PDF.}

% [review 2026-09-27] fix: the recursive definition of S applies for n > 2 (n = 2 is a base case), so the recurrence is for n >= 3 (n >= 2 would need N(0))
Here is the recurrence:  $N(1) = 0$, $N(2) = 1$, and, for any $n\ge 3$,

\begin{blue}
  \[
    N(n) = \REPLACEME
  \]
\end{blue}

\begin{lemma}
  $N(n) \le n\log_2 n$ for every $n\ge 1$.
\end{lemma}

\begin{longFormProof}

  \step The proof is by induction on $n$.

  \step\label{step: bus base 1}
  The base case $n=1$ holds because each set $S(\ell, \ell)$ has no segments, and $n\log_2 n$ is zero.
  
  \step\label{step: bus base 2}
  % [review 2026-09-27] fix: for n = 2, n log_2 n = 2, not one
  The base case $n=2$ holds because each set $S(\ell, \ell+1)$ has one segment, and $n\log_2 n$ is two.
  
  \begin{block}\label{block: bus induction}

    \step\label{step: bus ind assumption}
    Consider any $n\ge 3$ and assume the lemma holds for smaller values.  We show it holds for $n$.

    \step Recall the recurrence relation,  $N(n)$ is \BLUE{
      \(
      \REPLACEME
      \)
    }

    \step Applying the inductive assumption in Step~\ref{step: bus ind assumption}
    \BLUE{
      \[
        N(n)
        \le
        \REPLACEME
      \]
    }
    
    \step \BLUE{
      \REPLACEME
    }
  \end{block}

  \step By Block~\ref{block: bus induction}, for any $n\ge 3$, the lemma holds for that $n$ if it holds for smaller values.

  \step By this, Steps~\ref{step: bus base 1} and~\ref{step: bus base 2}, and induction, the lemma holds for all $n\ge 1$.

\end{longFormProof}

\continued 

\paragraph{Problem 3(d) answer}\GradescopeComment{This should be on page 12 (top) of your PDF.}

\begin{lemma}
  For any $n\ge 1$, for any $n$ consecutive stops $(\ell, \ell+1, \ldots, u)$,
  for every pair of stops $(i, j)$ with $\ell \le i < j \le u$,
  the hop length from $i$ to $j$ in $S(\ell, u)$ is at most 2.
\end{lemma}
\begin{longFormProof}
  \step The proof is by induction on $n$.

  \step\label{step: bus 2 base 1}
  For the base case $n=1$, there is no such pair $(i, j)$ so the lemma holds trivially. 

  \step\label{step: bus 2 base 2}
  For the base case $n=2$,
  there is one such pair $(\ell, \ell+1)$.
  \BLUE{
    \REPLACEME
  }

  \smallskip
  
  \begin{block}\label{block: bus 2 induction}
    \step\label{step: assumption}
    Consider any $n\ge 3$ and $(\ell, u)$ as in the lemma, and assume the lemma holds for smaller values.
    
    \step Let $m=\lfloor (\ell+u)/2 \rfloor$ be the middle stop.
    
    \step Recall that by definition $S(\ell, u)$ is \BLUE{
      \REPLACEME
    }
    
    \smallskip
    
    \begin{block}\label{block: bus ij}
      \step Consider any pair of stops $(i, j)$ with $\ell \le i < j \le u$.

      \begin{block}
        \step \underline{Case 1.} First consider the case that \BLUE{
          \(
          \REPLACEME
          \)
        }

        \step By the definition of $Q(\ell, u)$, the hop length from $i$ to $j$ is at most two in $Q(\ell, u)$.

        \step So the hop length from $i$ to $j$ is at most two in $S(\ell, u)$.
      \end{block}

      \begin{block}
        \step \underline{Case 2.} Next consider the case that \BLUE{
          \REPLACEME
        }

        \step \BLUE{
          \REPLACEME
        }

        \step So the hop length from $i$ to $j$ is at most two in $S(\ell, u)$. 
      \end{block}

      \begin{block}
        \step \underline{Case 3.} In the remaining case \BLUE{
          \REPLACEME
        }

        \step \BLUE{
          \REPLACEME
        }
        \step So the hop length from $i$ to $j$ is at most two in $S(\ell, u)$. 
      \end{block}

      \step In any case, the hop length from $i$ to $j$ in $S(\ell, u)$ is at most two.

    \end{block}

    \smallskip
    
    \step By Block~\ref{block: bus ij}, for all $(i, j)$ with $\ell\le i < j \le u$,
    the $i$-to-$j$ hop length in $S(\ell, u)$ is at most two.

  \end{block}

  \step By Block~\ref{block: bus 2 induction}, for any $n\ge 3$, the lemma holds for $n$ if it holds for smaller values.

  \step By this, Steps~\ref{step: bus 2 base 1} and~\ref{step: bus 2 base 2},
  and induction, the lemma holds for all $n\ge 1$.
\end{longFormProof}


\continued 


\paragraph{Problem \arabic{problem} collaboration and citations}~\GradescopeComment{This should be on page 13 (top) of your PDF.}

I collaborated on this problem with the following people:
\begin{blue}
  \begin{itemize}
  \item \REPLACEME          % write "none" if none
  \end{itemize}
\end{blue}

My answers use ideas from the following sources (other than the lecture notes and textbooks): 
\begin{blue}
  \begin{itemize}
  \item \REPLACEME         % write "none" if none
  \end{itemize}
\end{blue}

%%% Local Variables:
%%% mode: latex
%%% TeX-master: "homework_2"
%%% End:
